\documentclass[11pt]{article}
\usepackage[T1]{fontenc}
\usepackage[utf8]{inputenc}
\usepackage{amsmath,amsthm,mathtools,aliascnt}
\usepackage{newtxtext,newtxmath}
\usepackage[a4paper,margin=27mm,headheight=15pt]{geometry}
\usepackage{microtype}
\usepackage{enumitem,booktabs,longtable,array}
\usepackage{xcolor}
\usepackage{fancyhdr}
\usepackage{tikz}
\usetikzlibrary{arrows.meta,positioning}
\usepackage[most]{tcolorbox}
\usepackage{hyperref}
\usepackage[nameinlink,noabbrev]{cleveref}
\definecolor{ink}{HTML}{17324D}
\definecolor{soft}{HTML}{F2F5F8}
\hypersetup{colorlinks=true,linkcolor=ink,citecolor=ink,urlcolor=ink,
 pdftitle={Discrete Initial Groups and Omnific Normalization},
 pdfauthor={AI-assisted research draft prepared for the Surreal project},
 pdfsubject={Sign-tree surgery and a proposed solution of an Ehrlich--Kaplan question}}
\pagestyle{fancy}
\fancyhf{}
\fancyhead[L]{\small\textsc{Discrete initial groups}}
\fancyhead[R]{\small\textsc{Omnific normalization}}
\fancyfoot[C]{\thepage}
\renewcommand{\headrulewidth}{0.3pt}
\setlength{\parskip}{3pt}
\setlength{\emergencystretch}{2em}
\setcounter{tocdepth}{2}
\setlist[enumerate]{leftmargin=*,itemsep=3pt,topsep=4pt}
\newtheorem{theorem}{Theorem}[section]
% Shared counter through aliascnt, so that cleveref names each type
% correctly (with a plain [theorem] counter every reference read "theorem").
\newcommand{\isgthm}[3]{%
 \newaliascnt{#1}{theorem}%
 \newtheorem{#1}[#1]{#2}%
 \aliascntresetthe{#1}%
 \crefname{#1}{#3}{#3s}%
 \Crefname{#1}{#2}{#2s}}
\isgthm{lemma}{Lemma}{lemma}
\isgthm{proposition}{Proposition}{proposition}
\newaliascnt{corollary}{theorem}
\newtheorem{corollary}[corollary]{Corollary}
\aliascntresetthe{corollary}
\crefname{corollary}{corollary}{corollaries}
\Crefname{corollary}{Corollary}{Corollaries}
\isgthm{fact}{Imported fact}{imported fact}
\theoremstyle{definition}
\isgthm{definition}{Definition}{definition}
\isgthm{example}{Example}{example}
\theoremstyle{remark}
\isgthm{remark}{Remark}{remark}
\isgthm{question}{Question}{question}
\newtheorem{maintheorem}{Theorem}
\renewcommand{\themaintheorem}{\Alph{maintheorem}}
\crefname{maintheorem}{Theorem}{Theorems}
\newcommand{\No}{\ensuremath{\mathbf{No}}}
\newcommand{\Oz}{\ensuremath{\mathbf{Oz}}}
\newcommand{\On}{\ensuremath{\mathbf{On}}}
\newcommand{\R}{\mathbb{R}}
\newcommand{\C}{\mathbb{C}}
\newcommand{\Z}{\mathbb{Z}}
\newcommand{\N}{\mathbb{N}}
\newcommand{\D}{\mathbb{D}}
\newcommand{\supp}{\operatorname{supp}}
\newcommand{\Pred}{\operatorname{Pred}}
\newcommand{\Rs}{\operatorname{R}_{s}}
\newcommand{\Ls}{\operatorname{L}_{s}}
\newcommand{\len}{\ell}
\newcommand{\cat}{\mathbin{\frown}}
\newcommand{\plus}{\langle+\rangle}
\newcommand{\minus}{\langle-\rangle}
\newcommand{\emp}{\varnothing}
\newcommand{\pr}{\mathrel{\prec}}
\newcommand{\pre}{\mathrel{\preccurlyeq}}
\newcommand{\Hahn}[2]{\mathcal{H}_{#1}(#2)}
\newcommand{\Rect}{\mathcal{N}_{\alpha,n}}
\newcommand{\tht}{\Theta_{\alpha}}
\newcommand{\psiA}{\Psi_{\alpha}}
\newcommand{\ct}{\operatorname{ct}}
\newcommand{\lc}{\operatorname{lc}}
\newcommand{\hdeg}{\operatorname{deg}_{H}}
\newcommand{\lex}{\mathbin{\overrightarrow{\times}}}
\newcommand{\source}[1]{\textup{\cite{#1}}}
\newcommand{\replabel}[1]{\nolinkurl{#1}}
\newcommand{\fn}[1]{\nolinkurl{#1}}
\newtcolorbox{statusbox}{colback=soft,colframe=ink,boxrule=0.5pt,
 arc=1mm,left=3mm,right=3mm,top=2mm,bottom=2mm}

\begin{document}
\hypersetup{pageanchor=false}
\begin{titlepage}
\vspace*{2mm}
{\color{ink}\Large\textsc{Surreal numbers / ordered algebra}}\par
\vspace{6mm}
{\color{ink}\huge\bfseries Discrete Initial Groups\\[3mm]
and Omnific Normalization\par}
\vspace{5mm}
{\Large Sign-tree surgery, a sharp image classification,\\
and an Ehrlich--Kaplan question\par}
\vspace{7mm}
{\large AI-assisted research draft prepared for the Surreal project\par}
\vspace{2mm}
23 September 2026 \textperiodcentered{} single-source report of the
collection
\vspace{5mm}
\begin{abstract}
We give an explicit construction proposed as an affirmative solution of
Ehrlich and Kaplan's question whether every discretely ordered initial
subgroup of the surreal numbers is order-isomorphic to an initial subgroup
of the omnific integers (Question~2 of the arXiv preprint, Question~9.1 of
the journal version). The construction is not ordinary multiplication
by the reciprocal of the least positive element. Instead, a sign-sequence
operation moves the minimum exponent to zero, while preserving all
right-ancestor relations away from that minimum; a separate rescaling
changes only the bottom coefficient group to the ordinary integers.
The Ehrlich--Kaplan characterization of initial Hahn subgroups then
establishes initiality of the image.

The argument yields more than existence. For a prescribed least positive
element $2^{-n}\omega^{-\alpha}$, we characterize the images exactly by
containment of an explicit dyadic-spine subgroup $K_\alpha\subseteq\Oz$.
We identify the smallest and largest initial groups with that least
positive element, give sharp cardinality consequences, construct an
initial realization of the quotient by the least positive cyclic subgroup,
and extend the additive normalization to rectangular surcomplex modules.
All normal-form transports preserve leading truncations and set-indexed
strong summability. Counterexamples explain why scalar normalization,
exponent translation, and unrestricted inversion do not suffice.
\end{abstract}
\vspace{3mm}
\begin{statusbox}
\textbf{Status and scope.} This is an unrefereed research manuscript with
full proposed proofs, not an independently certified resolution or a
claim of established priority. The principal external mathematical input
is explicitly identified. Finite regression checks accompany the text;
they do not certify the transfinite argument. No new Lean formalization
is claimed. A targeted literature and repository review found the cited
question and its motivating example, but did not establish exhaustive
novelty. At placement in the collection every step was checked by hand
and no gap was found (\cref{isg:sec:provenance}); that check is not a
refereeing, and the literature search was targeted, not exhaustive.
\end{statusbox}
\vfill
{\small\textbf{Keywords:} surreal numbers; omnific integers; initial
subgroups; sign sequences; Hahn series; ordered abelian groups.\par
\textbf{Mathematical focus:} additive and simplicity-theoretic structure,
not multiplicative or exponential equivalence.}
\end{titlepage}
\hypersetup{pageanchor=true}

\begingroup
\small
\setlength{\parskip}{0pt}
\tableofcontents
\endgroup
\newpage

\section{The problem and the main results}\label{isg:sec:intro}

\subsection{A specific published question}
An additive subgroup of $\No$ is \emph{initial} when it contains every
proper sign-sequence prefix of each of its members. This is a much stronger
requirement than being an ordered subgroup. An initial subgroup is
\emph{discrete} when it is nonzero and has a least positive element.
Throughout, ``isomorphism'' of ordered groups preserves addition and
order; it need not preserve the ambient simplicity relation.

Ehrlich and Kaplan ask:
\begin{quote}
``Is every discrete initial subgroup of $\No$ isomorphic to an initial
subgroup of $\Oz$?''
\end{quote}
This is Question~2 in Section~9 of the explicitly versioned preprint
\cite{EK}, printed page~18. The journal article appeared in 2018
\cite{EKjournal}. We use the preprint's numbering throughout to avoid
mixing versions. The question also appears in Kaplan's thesis
\cite{Kaplan}.

In the journal version the same question is Question~9.1, with the same
wording and the same motivating example. The results of \cite{EK} used
below are numbered differently there; the correspondence, checked when
this report was placed in the collection against the author-hosted text
of the journal version \cite{EKjournal}, is:
\begin{center}\small
\begin{tabular}{@{}lll@{}}
\toprule
arXiv v1 \cite{EK} & journal version \cite{EKjournal} & use here\\
\midrule
Lemma~3 & Lemma~5.1 & initial subgroups of $\R$, \cref{isg:fact:real}\\
Theorem~1 & Theorem~5.1 & the initiality criterion, \cref{isg:fact:EK}\\
Proposition~9 & Proposition~5.2 & shape of the least positive element\\
Question~2 & Question~9.1 & the question addressed here\\
Question~3 & Question~9.2 & the initial-monoid problem, not addressed\\
\bottomrule
\end{tabular}
\end{center}
Below, the arXiv numbers are used, except where the wording of the
journal version matters (after \cref{isg:fact:EK}).

The distinction between a group and a ring matters. A discrete initial
subgroup can contain $1/2$ or a nonzero infinitesimal, whereas an omnific
integer has neither a nonintegral constant coefficient nor a negative
normal-form exponent. The task is therefore to change the initial
\emph{realization} of the ordered group, not to prove that its original
realization already lies in $\Oz$.

The paper addresses precisely this existence question and several
structural refinements. It does not claim a classification of all ordered
abelian groups admitting initial embeddings, nor a solution of the
separate initial-monoid problem (Question~3 of \cite{EK}, Question~9.2
of the journal version).

\subsection{Normalization and its exact range}
Write $\N=\{0,1,2,\ldots\}$ and $\D=\Z[1/2]$. If $\alpha$ is an ordinal, let
\[
 p_\alpha=-\alpha=\minus^{\alpha},
 \qquad q_\beta=\plus\cat\minus^{\beta}\quad(\beta\in\On).
\]
Here $\cat$ means concatenation of sign sequences, not surreal addition.
For example, $q_0=1$, $q_1=1/2$, and $q_m=2^{-m}$ for finite $m$.
Define the finite-support additive subgroup
\begin{equation}\label{isg:eq:Kintro}
 K_\alpha=\Z\oplus\bigoplus_{\beta<\alpha}\D\,\omega^{q_\beta}
 \ \subseteq\Oz.
\end{equation}
The direct sum means that each element has only finitely many nonzero
summands. In particular, $K_0=\Z$.

\begin{maintheorem}[Omnific normalization]\label{isg:thm:A}
Let $A\subseteq\No$ be a nonzero discrete initial additive subgroup.
There are unique $\alpha\in\On$ and $n\in\N$ such that its least positive
element is
\[
 \varepsilon=2^{-n}\omega^{-\alpha}.
\]
An explicitly defined map $\Rect$ restricts to an order-preserving group
isomorphism
\[
 \Rect:A\longrightarrow H,
 \qquad H\subseteq\Oz\text{ initial},\qquad \Rect(\varepsilon)=1.
\]
The map extends to a real-linear order-isomorphism between the full
normal-form spaces supported in $[-\alpha,\infty)$ and $[0,\infty)$.
This extension preserves support order types, all leading truncations,
and set-indexed strong summability in both directions.
\end{maintheorem}

\begin{maintheorem}[Sharp image classification]\label{isg:thm:B}
Fix $\alpha\in\On$ and $n\in\N$. Under the same explicit map $\Rect$,
the initial groups with least positive element
$2^{-n}\omega^{-\alpha}$ correspond bijectively, with inclusions
preserved and reflected, to the initial groups $H\subseteq\Oz$ satisfying
\[
 K_\alpha\subseteq H.
\]
The inverse is explicit. Thus the dyadic-spine condition is both necessary
and sufficient, not merely an artifact of a sufficient construction.
\end{maintheorem}

\begin{maintheorem}[Removing the bottom cyclic layer]\label{isg:thm:C}
For every group $A$ in \cref{isg:thm:A}, the subgroup $\Z\varepsilon$ is
convex, and the ordered quotient $A/\Z\varepsilon$ has an explicit
initial realization in $\No$. The corresponding short exact sequence
splits as an additive group, with a canonical splitting relative to the
given Conway normal forms. Equivalently,
\[
 A\cong B\lex\Z
\]
for an explicitly constructed initial subgroup $B\subseteq\No$, where
the first coordinate is dominant in the lexicographic order.
\end{maintheorem}

The main construction is in \cref{isg:sec:surgery,isg:sec:transport}; the proofs
of the three theorems are completed in
\cref{isg:sec:normalization,isg:sec:classification,isg:sec:quotient}, respectively.
Theorem~B contains more information than Theorem~A: it identifies exactly
which normalized realizations can be moved to any prescribed least
positive scale by this construction.

\subsection{Why a change of exponents is necessary}
Multiplication by $\varepsilon^{-1}$ certainly gives an ordered-group
isomorphism. Initiality, however, refers to sign prefixes rather than
numerical intervals, and multiplication does not generally preserve it.
Even $A=\frac12\Z+\Z\omega$ is an initial group for which
$2A=\Z+2\Z\omega$ is not initial: $2\omega$ belongs to the image but
its simpler predecessor $\omega$ does not.

Our operation treats two different issues separately. First, it relocates
the minimum of the \emph{exponent tree} to zero. Second, it rescales only
the real coefficient at that minimum. No claim of multiplicativity is
made, and none is needed for the published ordered-group question.

\subsection{Conventions shared with the collection}\label{isg:sec:conventions}
This report is built from one manuscript; its provenance, the placement
review and its relation to other reports are in \cref{isg:sec:provenance}.
Three conventions are fixed here once.
\begin{enumerate}
\item \emph{Initial} always means closed under proper sign-sequence
      prefixes, as in the definition of the definable-surreals report
      (\replabel{dsn:sec:foundations}). It never means an initial segment
      of the numerical order, and it has nothing to do with the
      \emph{initial forms} (leading forms of polynomial solutions) of the
      Diophantine report (\replabel{odg:prop:initial}).
\item $\Oz=\Z\oplus\Pi$, with $\Pi$ the class of purely infinite normal
      forms, is the omnific ring of the collection
      (\replabel{odg:def:rings}, \replabel{osq:eq:Ozdef},
      \replabel{found:eq:omnific}); \eqref{isg:eq:Oz} below is the same
      description. For $x\in\Oz$, $\ct(x)$ is its coefficient at
      exponent~$0$, as there. Only the additive and order structure of
      $\Oz$ is used; its ring structure plays no role.
\item The symbols $\Theta_\alpha$, $\Psi_\alpha$, $q_\beta$ and
      $\mathcal N_{\alpha,n}$ are local to this report. They are unrelated
      to the theta series and to the nome or dilation parameter $q$ of
      other reports.
\end{enumerate}

\section{Foundations and the imported initiality criterion}\label{isg:sec:background}

\subsection{Sets, classes, signs, and two orders}
We work in NBG with global choice, in the conventional class-sized
formulation of surreal numbers. Every individual surreal is a
set-length sign sequence. Every individual normal form has set support.
All strongly summable families considered below are indexed by sets.
Statements about families of class-sized groups are understood as
correspondence statements about class parameters, not as assertions that
a class of all classes exists.

For sign sequences $u,v$, write $u\pr v$ for a proper prefix and
$u\pre v$ for a prefix allowing equality. Their ordinary numerical order
is determined at the first disagreement, with
\[
 -\ <\ \text{undefined}\ <\ +.
\]
The empty sequence is the number $0$. Put
\[
 \Pred(x)=\{y:y\pr x\},\qquad
 \Rs(x)=\{y\pr x:x<y\},\qquad
 \Ls(x)=\{y\pr x:y<x\}.
\]
In sign language, a prefix obtained by stopping just before a minus sign
is a right ancestor; stopping just before a plus sign gives a left
ancestor. Each of these ancestor collections is a set.

A subclass $E\subseteq\No$ is initial precisely when
$\Pred(x)\subseteq E$ for every $x\in E$. This does not mean that $E$
is downward closed in the numerical order. For instance, $\No_{\ge0}$
is initial as a sign tree but is a numerical final segment.

If $u$ and $v$ have lengths $\lambda$ and $\mu$, then
$u\cat v$ has length $\lambda+_{\mathrm{ord}}\mu$. All additions of
\emph{lengths} in this paper are ordinal additions. Concatenation is
associative and preserves comparison between tails sharing the same
prefix. At a limit ordinal there is no last sign to remove: ``remove the
first sign'' or ``remove a specified concatenated block'' always means
restricting to the appropriate interval and taking its order type.
These conventions agree with the sign-sequence framework of
\cite{BagHoeven}.

\subsection{Hahn spaces without an exponent group}
For an ordered subclass $E\subseteq\No$, denote by $\Hahn{\R}{E}$
the additive space of formal expressions
\begin{equation}\label{isg:eq:nf}
 x=\sum_{\gamma\in S} r_\gamma\omega^\gamma,
 \qquad r_\gamma\in\R\setminus\{0\},
\end{equation}
where $S\subseteq E$ is a set reverse-well-ordered by $<$.
Conway normal forms identify these expressions with the corresponding
surreal numbers. The ordering is the sign of the coefficient at the
greatest exponent of a nonzero difference. For $x\ne0$, write
\[
 \hdeg(x)=\max\supp(x),\qquad
 \lc(x)=r_{\hdeg(x)}.
\]
We do not assert that $E$ is an additive group. Accordingly,
$\Hahn{\R}{E}$ is used here as an additive, real-linear Hahn space, not
necessarily as a ring. This distinction permits nonadditive exponent
reindexings.

Suppose $S=\{\gamma_\xi:\xi<\lambda\}$ is enumerated in strictly
decreasing order. A leading truncation of $x$ is
\[
 x_{<\eta}=\sum_{\xi<\eta}r_{\gamma_\xi}\omega^{\gamma_\xi}
 \quad(\eta\le\lambda).
\]
``Truncation closed'' will always refer to these leading truncations,
including those at limit indices, not to closure under arbitrary
subseries. The normal-form and monomial facts used here are part of the
standard theory summarized in \cite{EK}; we never choose a noncanonical
Hahn embedding of an arbitrary ordered group.

\begin{definition}[Strong summability]\label{isg:def:strong}
A set-indexed family $(x_i)_{i\in I}$ in a Hahn space is strongly
summable if $\bigcup_i\supp(x_i)$ is reverse-well-ordered and each
exponent occurs in only finitely many supports. Its sum is the normal
form obtained by the finite coefficient sum at each exponent, with zero
coefficients discarded.
\end{definition}

For a subgroup $G\subseteq\Hahn{\R}{E}$, set
\[
 G_\gamma=\{r\in\R:r\omega^\gamma\in G\}.
\]
It is \emph{cross-sectional on $E$} when $\omega^\gamma\in G$ for
every $\gamma\in E$.

\begin{fact}[Ehrlich--Kaplan, concrete form]\label{isg:fact:EK}
An initial subgroup of $\No$ is a truncation-closed, cross-sectional
Hahn subgroup on its initial exponent class
$\Gamma(G)=\{\gamma:\omega^\gamma\in G\}$. Its nonzero coefficient
groups are initial subgroups of $\R$, and
\begin{equation}\label{isg:eq:dyadic}
 \gamma,\eta\in\Gamma(G),\quad \eta\in\Rs(\gamma)
 \quad\Longrightarrow\quad \D\subseteq G_\eta.
\end{equation}
Conversely, these conditions imply that the canonical normal-form image
is initial in $\No$.
\end{fact}

This is the form of \cite[Theorem~1 and its proof]{EK} needed here.
The last sentence is important: the sufficiency proof establishes
initiality of the \emph{specified canonical image}, not just existence of
some other initial embedding.

The journal version states this concrete form outright. Its Theorem~5.1,
the counterpart of \cite[Theorem~1]{EK}, says that a subgroup of $\No$ is
initial if and only if it is isomorphic ``via the canonical isomorphism''
to a truncation-closed, cross-sectional Hahn subgroup whose exponent
class is an initial ordered subclass of $\No$ and whose coefficient
groups satisfy the two conditions above \cite{EKjournal}. The canonical
isomorphism is the normal-form identification of $\No$ with the Hahn
group, so the journal statement is \cref{isg:fact:EK} as used here.
The criterion is imported as a published statement. Neither the source
manuscript nor the placement review re-proves it
(\cref{isg:sec:provenance}).

\begin{fact}[Real coefficient groups]\label{isg:fact:real}
A nonzero initial additive subgroup of $\R$ either contains $\D$ or is
$2^{-m}\Z$ for an integer $m\ge0$ \cite[Lemma~3]{EK}
\textup{(}Lemma~5.1 of the journal version\textup{)}.
In particular, a discrete one has least positive element $2^{-m}$.
\end{fact}

No proof in the present paper replaces \cref{isg:fact:EK} by the weaker
assertion that initiality of the exponent set alone is sufficient. The
coefficient groups and the right-ancestor condition are essential.

\subsection{Omnific integers and coefficient isolation}
We use the normal-form description
\begin{equation}\label{isg:eq:Oz}
 \Oz=\left\{\sum_{\gamma>0}r_\gamma\omega^\gamma+m:
           m\in\Z\right\},
\end{equation}
with set reverse-well-ordered positive support and arbitrary real
coefficients. This is Conway's canonical integer part, as recalled in
\cite[Section~6.2]{Kaplan}, and the collection's $\Oz=\Z\oplus\Pi$
(\cref{isg:sec:conventions}). It is an initial additive subgroup: apply
\cref{isg:fact:EK} to $E=\No_{\ge0}$, coefficient group $\Z$ at zero and
$\R$ at every positive exponent. Zero is never a right ancestor of a
nonnegative number, so the exceptional coefficient group causes no
violation of \eqref{isg:eq:dyadic}. Every positive element of \Oz is at
least $1$; this discreteness is also \replabel{odg:prop:units} of the
Diophantine report.

\begin{lemma}[Isolating a normal-form term]\label{isg:lem:isolate}
Let $G$ be an additive subgroup closed under leading truncations. If
$x\in G$ has a nonzero coefficient $r_\gamma$ at exponent $\gamma$,
then $r_\gamma\omega^\gamma\in G$.
\end{lemma}
\begin{proof}
If $\gamma=\gamma_\eta$ in the decreasing enumeration of the support,
subtract the truncation at $\eta$ from the truncation at $\eta+1$.
Both belong to $G$. This also works when $\eta$ is a limit ordinal.
\end{proof}

This lemma permits removing a specified individual term. It does not
license an arbitrary infinite subseries of an element of $G$.

\section{Locating the least positive scale}\label{isg:sec:bottom}

\begin{lemma}[A minimum in an initial tree]\label{isg:lem:min-tree}
If a nonempty initial subclass $E\subseteq\No$ has a numerical minimum
$p$, then $p=-\alpha$ for a unique ordinal $\alpha$.
Moreover, the cone $C_\alpha=\{x\in\No:x\ge-\alpha\}$ is initial.
\end{lemma}
\begin{proof}
A plus sign in $p$ would give a proper prefix smaller than $p$.
Initiality would place that prefix in $E$, contradicting minimality.
Thus every sign in $p$ is minus.

For the second assertion, let $p=-\alpha$, $x\ge p$, and $v\pr x$.
If $v<p$, then at the first disagreement with the all-minus sequence $p$
one would have to find either a sign smaller than minus, which is
impossible, or a minus sign after $p$ has ended. The latter would force
$x$ itself to be below $p$. Hence $v\ge p$.
\end{proof}

\begin{proposition}[Bottom coefficient structure]\label{isg:prop:bottom}
Let $A$ be a nonzero discrete initial subgroup with least positive
member $\varepsilon$. There are unique $\alpha\in\On$, $n\in\N$
such that, writing $p=-\alpha$ and $c=2^{-n}$,
\begin{equation}\label{isg:eq:bottomdata}
 \varepsilon=c\omega^p,\qquad
 \min\Gamma(A)=p,\qquad A_p=c\Z.
\end{equation}
Every support of every member of $A$ lies in $C_\alpha$.
\end{proposition}
\begin{proof}
By \cref{isg:fact:EK}, $A$ is truncation closed. Write the leading term of
$\varepsilon$ as $r\omega^p$ with $r>0$. This term belongs to $A$.
If the remaining tail $t$ were nonzero, then $|t|\in A$ and
$0<|t|<\varepsilon$, since the leading exponent of $t$ is strictly
smaller than $p$. This contradicts the choice of $\varepsilon$.
Consequently $\varepsilon=r\omega^p$.

The exponent $p$ lies in $\Gamma(A)$ by the cross-sectional normal-form
description. If $\gamma<p$ belonged to $\Gamma(A)$, then
$0<\omega^\gamma<r\omega^p=\varepsilon$, another contradiction.
Thus $p$ is the minimum of $\Gamma(A)$. By \cref{isg:lem:min-tree},
$p=-\alpha$ for a unique ordinal $\alpha$.

The real coefficient group $A_p$ has least positive element $r$.
By \cref{isg:fact:real}, it is $2^{-n}\Z$ for a unique $n\in\N$, and
$r=2^{-n}$. All supports lie in $\Gamma(A)$, proving the last assertion.
\end{proof}

The shape of the least positive element is consistent with
\cite[Proposition~9 and its proof]{EK} (Proposition~5.2 of the journal
version); the argument above also records
the support and coefficient data needed later.

\begin{corollary}[Forced ancestors]\label{isg:cor:forced}
Under \eqref{isg:eq:bottomdata}, for every $\beta<\alpha$ one has
\[
 -\beta\in\Gamma(A),\qquad \D\omega^{-\beta}\subseteq A.
\]
\end{corollary}
\begin{proof}
The number $-\beta$ is a right ancestor of $-\alpha$.
Initiality gives the exponent, and \eqref{isg:eq:dyadic} gives all dyadic
coefficients there.
\end{proof}

The bottom coefficient group can be discrete, but all these ancestors
must carry dyadically divisible coefficient groups. This is the information
that survives normalization as the subgroup $K_\alpha$.

\section{Moving a minimum to zero in the sign tree}\label{isg:sec:surgery}

\subsection{Definition and inverse}
Fix $\alpha\in\On$, put $p=-\alpha$, and keep
$C_\alpha=[p,\infty)$. Every $x\in C_\alpha$ belongs to exactly one of
the following cases: $x=p$; $x=p\cat\plus\cat u$ for a unique tail
$u$; or $p$ is not a prefix of $x$.
Indeed, an extension through $p\cat\minus$ would lie below the cone.

\begin{definition}[Root relocation]\label{isg:def:theta}
Define $\tht:C_\alpha\to\No_{\ge0}$ by
\begin{equation}\label{isg:eq:theta}
 \tht(x)=
 \begin{cases}
  0, & x=p,\\
  \plus\cat p\cat u, & x=p\cat\plus\cat u,\\
  \plus\cat x, & p\not\pre x.
 \end{cases}
\end{equation}
The inverse candidate $\psiA:\No_{\ge0}\to C_\alpha$ is
\begin{equation}\label{isg:eq:psi}
 \psiA(0)=p,\qquad
 \psiA(\plus\cat v)=
 \begin{cases}
  p\cat\plus\cat u,&v=p\cat u,\\
  v,&p\not\pre v.
 \end{cases}
\end{equation}
\end{definition}

The symbols in these formulas are sign blocks. For $\alpha=0$, the block
$p$ is empty and both maps are the identity. For $\alpha>0$, one may
picture the operation as deleting the old minimum node, moving its
right subtree into the vacated position, and adding a new common root.
The formulas, not this picture, govern limit-length cases.

\begin{lemma}[The inverse is well-defined]\label{isg:lem:inverse}
The maps \eqref{isg:eq:theta} and \eqref{isg:eq:psi} are inverse bijections
between $C_\alpha$ and $\No_{\ge0}$.
\end{lemma}
\begin{proof}
Every nonempty image under $\tht$ begins with plus and is therefore
positive. Conversely, let $y=\plus\cat v>0$. If $p\pre v$, the
specified preimage extends $p$ through plus and is greater than $p$.
If $p\not\pre v$, then $v>p$: a sequence below an all-minus word must
extend that entire word through minus, contrary to the assumption.
Thus $\psiA$ takes values in the cone.

If $x=p$, the composite $\psiA\tht$ returns $p$. If
$x=p\cat\plus\cat u$, removing the first sign of
$\plus\cat p\cat u$ leaves $p\cat u$, and the first inverse case
returns $x$. In the remaining case, removing the first sign leaves
$x$, to which the second inverse case applies. The same two cases for
$v$ prove $\tht\psiA(y)=y$ for positive $y$, and zero is immediate.
Uniqueness of tails follows from the order types of the corresponding
ordinal intervals, so none of these steps requires a successor length.
\end{proof}

\begin{lemma}[Order preservation]\label{isg:lem:thetaorder}
The bijection $\tht$ is strictly increasing.
\end{lemma}
\begin{proof}
The minimum $p$ maps to the minimum $0$. On the block
\[
 U=\{p\cat\plus\cat u:u\in\No\},
\]
comparison is comparison of the tails $u$, before and after the map.
On the remaining class $V=C_\alpha\setminus(U\cup\{p\})$, adding
one common initial plus sign also preserves comparison.

It remains to compare the two blocks. Every member of $U$ is smaller
than every member $v$ of $V$: before the end of $p$, such a $v$ either
stops or first has plus where $p$ has minus. In both cases $v$ is greater
than every extension of $p$. After applying $\tht$, cancel the common
initial plus sign. One then compares an extension $p\cat u$ of $p$
with $v$, and the same first disagreement proves the same inequality.
When $\alpha=0$, the block $V$ is empty. This covers all cases.
\end{proof}

\subsection{The exact predecessor identity}
The following identity is the central combinatorial statement.
It makes the exception at the old minimum completely explicit.

\begin{theorem}[Predecessor and right-ancestor transport]\label{isg:thm:tree}
For every $x\in C_\alpha$ with $x>p$,
\begin{equation}\label{isg:eq:predidentity}
 \Pred(\tht(x))=\tht\bigl[\Pred(x)\cup\{p\}\bigr].
\end{equation}
Consequently,
\begin{align}
 \Rs(\tht(x))&=\tht[\Rs(x)],\label{isg:eq:rightidentity}\\
 \Ls(\tht(x))&=\tht[\Ls(x)\cup\{p\}].\label{isg:eq:leftidentity}
\end{align}
At $x=p$ the image has no proper ancestors, while
\begin{equation}\label{isg:eq:lost}
 \Rs(p)=\Pred(p)=\{-\beta:\beta<\alpha\},\qquad
 \tht[\Rs(p)]=\{q_\beta:\beta<\alpha\}.
\end{equation}
For $x,y>p$, simplicity is preserved and reflected:
\[
 x\pr y\quad\Longleftrightarrow\quad\tht(x)\pr\tht(y).
\]
\end{theorem}
\begin{proof}
First suppose $p\not\pre x$. No prefix $v\pre x$ can equal or extend
$p$; otherwise $p\pre x$. Thus $\tht(v)=\plus\cat v$ for all
proper prefixes $v$ of $x$. The proper prefixes of $\plus\cat x$ are
exactly the empty word, together with these $\plus\cat v$.
The empty word is $\tht(p)$, proving \eqref{isg:eq:predidentity} in this
case.

Now suppose $x=p\cat\plus\cat u$. The proper prefixes of $x$ are
\[
 \{-\beta:\beta\le\alpha\}
 \ \cup\ \{p\cat\plus\cat v:v\pr u\}.
\]
Under $\tht$ these become, respectively,
\[
 \{q_\beta:\beta<\alpha\}\cup\{0\},
 \qquad \{\plus\cat p\cat v:v\pr u\}.
\]
These are precisely the proper prefixes of $\plus\cat p\cat u$.
When $u$ is empty the second collection is empty; when $u$ has limit
length it includes exactly its proper ordinal restrictions. This proves
\eqref{isg:eq:predidentity} without a missing limit case.

Filter \eqref{isg:eq:predidentity} according to numerical comparison with
$\tht(x)$. By \cref{isg:lem:thetaorder}, comparison of images is comparison
of preimages. Since $p<x$, the extra element $\tht(p)=0$ is a left
ancestor, never a right ancestor. This gives
\eqref{isg:eq:rightidentity} and \eqref{isg:eq:leftidentity}.
Formula \eqref{isg:eq:lost} follows directly from the all-minus word $p$.
Finally, if $x,y>p$, membership of $\tht(x)$ in the predecessor set of
$\tht(y)$ is, by injectivity and \eqref{isg:eq:predidentity}, exactly
membership of $x$ in $\Pred(y)$.
\end{proof}

\begin{remark}[A stronger identity that would be false]\label{isg:rem:exception}
For $\alpha>0$ one must not assert \eqref{isg:eq:rightidentity} at $x=p$.
Already $p=-1$ has right ancestor $0$, whose image is $1$, while
$\tht(p)=0$ has no ancestors at all. The missing relations at precisely
this one point account for the extra condition in \cref{isg:thm:B}.
\end{remark}

\begin{corollary}[Initiality of the image]\label{isg:cor:initialimage}
If $\Gamma$ is initial and has minimum $p=-\alpha$, then
$\tht[\Gamma]$ is initial and has minimum $0$.
\end{corollary}
\begin{proof}
The minimum assertion follows from the order-isomorphism. The image of
$p$ has no predecessors. For $x>p$ in $\Gamma$, the right side of
\eqref{isg:eq:predidentity} lies in $\tht[\Gamma]$, because both
$\Pred(x)$ and $p$ lie in $\Gamma$.
\end{proof}

\begin{proposition}[The exact inverse obstruction]\label{isg:prop:inversetree}
Let $\Delta\subseteq\No_{\ge0}$ be nonempty and initial. Put
\[
 \Lambda_\alpha=\{0\}\cup\{q_\beta:\beta<\alpha\}.
\]
Then $\psiA[\Delta]$ is initial if and only if
$\Lambda_\alpha\subseteq\Delta$.
\end{proposition}
\begin{proof}
The empty word $0$ belongs to $\Delta$, so $p$ belongs to the inverse
image. If the inverse image is initial, it contains every $-\beta$ with
$\beta<\alpha$; applying $\tht$ gives $q_\beta\in\Delta$.

Conversely, suppose $\Lambda_\alpha\subseteq\Delta$. All predecessors
of $p$ are then in the inverse image. For any other inverse-image element
$x=\psiA(y)$, where $y>0$, \eqref{isg:eq:predidentity} gives
\[
 \Pred(x)\cup\{p\}=\psiA[\Pred(y)]\subseteq\psiA[\Delta].
\]
Hence every inverse-image element has all its predecessors in that image.
\end{proof}

\subsection{Exponent lengths and a finite picture}
\begin{proposition}[An exponent-birthday bound]\label{isg:prop:length}
For every $x\in C_\alpha$,
\[
 \len(\tht(x))\le 1+_{\mathrm{ord}}\len(x).
\]
In particular, an infinite exponent length never increases.
\end{proposition}
\begin{proof}
At $x=p$ the image length is zero. In the outside case the length is
exactly $1+\len(x)$. In the inside case, if $\len(u)=\mu$, the source
length is $\alpha+1+\mu$, and the image length is $1+\alpha+\mu$.
Monotonicity of ordinal addition gives
$1+\alpha+\mu\le1+\alpha+1+\mu$. Finally, $1+\lambda=\lambda$ for
any infinite ordinal $\lambda$: write $\lambda$ as a nonzero limit
ordinal followed by a finite tail.
\end{proof}

This is a bound for \emph{exponents}. It is not a theorem bounding the
birthday of the entire transformed surreal normal form.

\begin{center}
\begin{tikzpicture}[>=Stealth,every node/.style={font=\small},
                    x=1.30cm,y=1.0cm]
\node (a) at (0,0) {$0$};
\node (b) at (-1,-1) {$-1$};
\node (c) at (1,-1) {$1$};
\node (d) at (-0.5,-2) {$-\tfrac12$};
\node (e) at (0.6,-2) {$\tfrac12$};
\draw (a)--(b);\draw(a)--(c);\draw(b)--(d);\draw(c)--(e);
\node at (0,-2.65) {Source: a finite part of $C_1$};
\node (r) at (5,0) {$0$};
\node (s) at (5,-1) {$1$};
\node (t) at (6,-2) {$2$};
\node (u) at (4,-2) {$\tfrac12$};
\node (v) at (5.6,-3) {$\tfrac32$};
\draw(r)--(s);\draw(s)--(t);\draw(s)--(u);\draw(t)--(v);
\draw[->,thick] (1.8,-0.7)--node[above]{$\Theta_1$}(3.2,-0.7);
\node at (5,-3.55) {Image and its prefix relations};
\end{tikzpicture}
\end{center}
Here $-1\mapsto0$, $0\mapsto1$, $1\mapsto2$,
$-\tfrac12\mapsto\tfrac12$, and
$\tfrac12\mapsto\tfrac32$. The numerical labels are exponents, not
the monomials with those exponents.

\section{Strong transport of normal forms}\label{isg:sec:transport}

\begin{definition}[Bottom-normalized transport]\label{isg:def:normalization}
For $\alpha\in\On$, $n\in\N$, $p=-\alpha$, and $c=2^{-n}$, define
\begin{equation}\label{isg:eq:N}
 \Rect\left(\sum_{\gamma\ge p}r_\gamma\omega^\gamma\right)
   =c^{-1}r_p+
     \sum_{\gamma>p}r_\gamma\omega^{\tht(\gamma)}.
\end{equation}
A coefficient not in the support is interpreted as zero. The inverse is
\begin{equation}\label{isg:eq:Ninv}
 \Rect^{-1}\left(s_0+\sum_{\delta>0}s_\delta\omega^\delta\right)
   =c s_0\omega^p+
     \sum_{\delta>0}s_\delta\omega^{\psiA(\delta)}.
\end{equation}
\end{definition}

\begin{theorem}[Ambient transport]\label{isg:thm:transport}
The formulas above define inverse real-linear order-isomorphisms
\[
 \Rect:\Hahn{\R}{C_\alpha}
       \longrightarrow\Hahn{\R}{\No_{\ge0}}.
\]
They preserve support cardinality and decreasing support order type,
commute with every leading truncation, and preserve and reflect strong
summability of set-indexed families. For every such strongly summable
family,
\begin{equation}\label{isg:eq:strongtransport}
 \Rect\left(\sum_{i\in I}x_i\right)
       =\sum_{i\in I}\Rect(x_i).
\end{equation}
For a nonzero $x$,
\begin{equation}\label{isg:eq:leadingtransport}
 \hdeg(\Rect(x))=\tht(\hdeg(x)),
 \qquad
 \lc(\Rect(x))=
 \begin{cases}
  c^{-1}\lc(x),&\hdeg(x)=p,\\
  \lc(x),&\hdeg(x)>p.
 \end{cases}
\end{equation}
\end{theorem}
\begin{proof}
An order-isomorphism takes a reverse-well-ordered support to a
reverse-well-ordered support and preserves its order type. The support is
still a set: it is the image of a set under a definable class function.
Every coefficient is multiplied by a nonzero real number, so no support
point is lost and no new point is created. These observations show that
both displayed formulas are legitimate normal forms. Their inverse
identities follow coefficient by coefficient from
\cref{isg:lem:inverse}.

Both addition of normal forms and the proposed map are coefficientwise,
with the same fixed multiplier $c^{-1}$ at the bottom point. Hence the
map is real-linear. The greatest nonzero coefficient retains its sign,
because all multipliers are positive. Thus the map and its inverse
preserve order. This proves \eqref{isg:eq:leadingtransport} as well.

A leading truncation is specified by an initial segment of the decreasing
support enumeration. The same enumeration indices specify the
corresponding truncation of the transported support. Applying the map
therefore commutes with truncation, including at a limit index.

For a set-family $(x_i)$, one has the exact equality
\[
 \bigcup_i\supp(\Rect(x_i))
 =\tht\left[\bigcup_i\supp(x_i)\right].
\]
Reverse well-ordering of these unions is equivalent in both directions.
For each $\gamma$, the index set of family members contributing at
$\gamma$ is exactly the index set contributing at $\tht(\gamma)$ after
transport. The point-finiteness requirement is therefore also equivalent.
When it holds, coefficient extraction reduces \eqref{isg:eq:strongtransport}
to a finite real sum, through which the fixed coefficient multiplier
commutes. This proves the entire assertion without invoking convergence
in an order topology.
\end{proof}

\begin{remark}[No implicit completion of a subgroup]\label{isg:rem:nocompletion}
The ambient Hahn spaces in \cref{isg:thm:transport} are full spaces. A
particular initial subgroup $A$ need not be closed under all strong sums
or all infinite subseries. We define its image as the actual set or class
$\Rect[A]$, not as the space of all series with the same coefficient
groups. If a strongly summable family in $A$ has sum outside $A$,
\eqref{isg:eq:strongtransport} still holds in the ambient spaces, but is not
an assertion of an internal summation operation on $A$.
\end{remark}

\begin{remark}[What is and is not preserved]\label{isg:rem:preserved}
The map preserves the ordered chain of Archimedean classes, with their
exponents reindexed by $\tht$. It is not a homomorphism of additive
\emph{exponent} groups. Consequently \eqref{isg:eq:leadingtransport} should
not be described as preservation of a multiplicative valuation with the
same value-group structure.
\end{remark}

\section{Proof of omnific normalization}\label{isg:sec:normalization}

We now prove \cref{isg:thm:A}. Let $A$ be as there, and obtain
$\alpha,n,p,c$ from \cref{isg:prop:bottom}. Write
\[
 \Gamma=\Gamma(A),\qquad \Delta=\tht[\Gamma],
 \qquad H=\Rect[A].
\]
The ambient transport restricts to an ordered-group isomorphism onto
$H$. It remains to prove that this actual image is an initial subgroup
of $\Oz$.

\begin{lemma}[Transported coefficients]\label{isg:lem:coefficients}
The group $H$ is truncation closed and cross-sectional on $\Delta$,
and its coefficient groups are
\begin{equation}\label{isg:eq:coefftransport}
 H_0=\Z,\qquad H_{\tht(\gamma)}=A_\gamma\quad(\gamma>p).
\end{equation}
Its exponent class is exactly $\Delta$.
\end{lemma}
\begin{proof}
Truncation closure follows from \cref{isg:thm:transport}. For
$\gamma>p$, one has
$\Rect(\omega^\gamma)=\omega^{\tht(\gamma)}$. At the bottom,
$\Rect(c\omega^p)=1$. Thus all the leaders on $\Delta$ belong to $H$.
No image element can have a support point outside $\Delta$.

The formula at a positive image exponent leaves its coefficient
unchanged, and the inverse formula has the same property. This proves
both inclusions in the second identity of \eqref{isg:eq:coefftransport}.
At zero the coefficient group is $c^{-1}A_p=\Z$.
\end{proof}

\begin{proof}[Proof of \cref{isg:thm:A}]
By \cref{isg:cor:initialimage}, $\Delta$ is initial with minimum zero.
By \cref{isg:lem:coefficients}, the image is truncation closed and
cross-sectional, and all its coefficient groups are initial real
subgroups.

Consider an obligation in the dyadic right-ancestor condition for $H$:
let $\delta,\eta\in\Delta$ and $\eta\in\Rs(\delta)$.
The point $\delta$ cannot be zero. Write $\delta=\tht(\gamma)$
with $\gamma>p$. By \eqref{isg:eq:rightidentity}, there exists
$\xi\in\Rs(\gamma)$ with $\eta=\tht(\xi)$. Since $\xi>\gamma>p$,
its coefficient group was not rescaled. The initiality of $A$ gives
\[
 \D\subseteq A_\xi=H_\eta.
\]
Every target obligation is therefore satisfied. The concrete sufficiency
in \cref{isg:fact:EK} proves that $H$ is initial in $\No$.

All exponents of $H$ are nonnegative and its zero coefficient is an
integer. Hence $H\subseteq\Oz$ by \eqref{isg:eq:Oz}. Its minimum positive
element is $1$, either by order-isomorphism or directly from its normal
forms. Finally, $\Rect(\varepsilon)=1$, and all support and summability
claims are the restrictions of \cref{isg:thm:transport}. Uniqueness of
$\alpha,n$ was established in \cref{isg:prop:bottom}.
\end{proof}

\begin{corollary}[The cited existence question]\label{isg:cor:answer}
Every discrete initial subgroup of $\No$ is isomorphic, as an ordered
abelian group, to an initial subgroup of $\Oz$.
\end{corollary}

\begin{remark}[An initial image is not a simplicity isomorphism]\label{isg:rem:notsimplicity}
The conclusion is exactly an ordered-group isomorphism to an initial
image. For example, on $A=\frac12\Z$ the map $\mathcal N_{0,1}$
sends $1$ to $2$ and $1/2$ to $1$. Although $1\pr1/2$, one does not
have $2\pr1$. Thus the normalization is not generally an isomorphism
of the full simplicity-enriched structures. Any initial cyclic subgroup
of $\Oz$ is $\Z$, and the order-isomorphism from $\frac12\Z$ to
$\Z$ is forced. Thus this example rules out a generally
simplicity-preserving version. The published question does not demand
that stronger conclusion.
\end{remark}

\section{Exact images, extremal groups, and admissible scales}\label{isg:sec:classification}

\subsection{The compulsory dyadic spine}
\begin{lemma}[The spine group is initial]\label{isg:lem:K}
For every ordinal $\alpha$, the group $K_\alpha$ of
\eqref{isg:eq:Kintro} is initial in $\No$, has least positive element $1$,
and lies in $\Oz$. Its exponent class is $\Lambda_\alpha$.
\end{lemma}
\begin{proof}
The proper prefixes of $q_\beta$ are $0$ and the $q_\eta$ with
$\eta<\beta$. Thus $\Lambda_\alpha$ is initial. The finite-support
group on this class is cross-sectional and truncation closed. Its
coefficient group is $\Z$ at zero and $\D$ at every other point.
Any right ancestor lies among those positive spine points, where the
required dyadic containment holds. Apply \cref{isg:fact:EK}. The remaining
claims follow from its normal forms.
\end{proof}

\begin{proposition}[Necessary image data]\label{isg:prop:Knecessary}
If $A$ has least positive element $2^{-n}\omega^{-\alpha}$ and is
initial, then $K_\alpha\subseteq\Rect[A]$.
\end{proposition}
\begin{proof}
The image contains $1$. For each $\beta<\alpha$, \cref{isg:cor:forced}
gives $\D\omega^{-\beta}\subseteq A$. These terms are above the
bottom exponent, so their coefficients are unchanged, and
$\tht(-\beta)=q_\beta$. Thus the image contains every generating
summand of $K_\alpha$ and hence all their finite sums.
\end{proof}

\begin{proof}[Proof of \cref{isg:thm:B}]
Necessity has just been proved. Conversely, let $H\subseteq\Oz$ be
initial with $K_\alpha\subseteq H$, and put $A=\Rect^{-1}[H]$.
Since $H$ contains $1$, it has least positive element $1$ and $H_0=\Z$.
Let $\Delta=\Gamma(H)$. The inclusion of $K_\alpha$ implies
$\Lambda_\alpha\subseteq\Delta$, so
$\Gamma=\psiA[\Delta]$ is initial by \cref{isg:prop:inversetree}.
Its minimum is $p$.

By ambient transport, $A$ is an additive, truncation-closed Hahn subgroup.
Its bottom coefficient group is $c\Z$, which contains $1$ since
$c=2^{-n}$; all other coefficient groups are inherited from $H$.
It is therefore cross-sectional and has initial real coefficient groups.

We check the source right-ancestor condition. If $x\in\Gamma$ satisfies
$x>p$ and $y\in\Rs(x)$, then
$\tht(y)\in\Rs(\tht(x))$. The coefficient group at $y$ is unchanged
from that at $\tht(y)$, so it contains $\D$ by the initiality of $H$.
If instead $x=p$, its right ancestors are exactly the $-\beta$ for
$\beta<\alpha$. The hypothesis $K_\alpha\subseteq H$ says that all
coefficient groups at their images $q_\beta$ contain $\D$; again these
coefficients are unchanged. This checks the only obligations not already
covered by \eqref{isg:eq:rightidentity}.

The criterion \cref{isg:fact:EK} now proves $A$ initial. Its least positive
element is $\Rect^{-1}(1)=c\omega^p$, since the inverse is increasing.
The inverse identities show uniqueness, and a fixed bijection preserves
and reflects inclusion of its images. This proves the correspondence.
\end{proof}

\subsection{Smallest and largest initial realizations at a fixed scale}
Define
\begin{align}
 L_{\alpha,n}
   &=c\Z\omega^p\ \oplus\!
     \bigoplus_{\beta<\alpha}\D\omega^{-\beta},\label{isg:eq:L}\\
 M_{\alpha,n}
   &=\{x\in\Hahn{\R}{C_\alpha}:[\omega^p]x\in c\Z\}.
   \label{isg:eq:M}
\end{align}
Here $[\omega^p]x$ denotes coefficient extraction. The direct sum in
\eqref{isg:eq:L} again permits only finite supports; the class in
\eqref{isg:eq:M} permits every set reverse-well-ordered support in the cone.

\begin{theorem}[Extremal groups]\label{isg:thm:extremal}
The groups $L_{\alpha,n}$ and $M_{\alpha,n}$ are initial with least
positive element $c\omega^p$. Every initial group $A$ with that least
positive element satisfies
\[
 L_{\alpha,n}\subseteq A\subseteq M_{\alpha,n}.
\]
Moreover,
\begin{equation}\label{isg:eq:extremalimages}
 \Rect[L_{\alpha,n}]=K_\alpha,
 \qquad \Rect[M_{\alpha,n}]=\Oz,
 \qquad M_{\alpha,n}=c\omega^p\Oz.
\end{equation}
\end{theorem}
\begin{proof}
The first two identities in \eqref{isg:eq:extremalimages} follow directly
from the formulas: every positive image coefficient is unrestricted in
$M$, whereas precisely the finite dyadic spine terms occur for $L$.
By \cref{isg:thm:B}, the inverse images of $K_\alpha$ and $\Oz$ are
initial and have the required minimum. The lower bound on $A$ follows
from \cref{isg:cor:forced} and $A_p=c\Z$; the upper bound follows from
\cref{isg:prop:bottom}.

For the last identity, ordinary multiplication by $c\omega^p$ sends a
normal form on $[0,\infty)$ to one on $[p,\infty)$ by the arithmetic
translation $\delta\mapsto p+\delta$. Its bottom coefficient belongs
to $c\Z$, and the other coefficients range over all of $\R$ because
multiplication by $c\ne0$ is a bijection of $\R$. Conversely, translate
any support in the cone back by $-p$ and divide its coefficients by $c$.
The resulting series belongs to $\Oz$.
\end{proof}

The equality $M_{\alpha,n}=\varepsilon\Oz$ does not make the theorem
trivial. It concerns the \emph{largest} group, whose higher coefficients
are all real and whose support class is the full cone. Arithmetic
translation can destroy initiality for smaller initial subgroups, even
though it behaves correctly for this maximal one.

\begin{corollary}[Sharp size obstruction]\label{isg:cor:size}
Every set-sized initial group $A$ with least positive element
$2^{-n}\omega^{-\alpha}$ satisfies
\[
 |A|\ge\max(\aleph_0,|\alpha|).
\]
For every ordinal $\alpha$ and every $n\in\N$ this bound is attained
by $L_{\alpha,n}$. In particular, a countable initial group can have
such a least positive element only when $\alpha$ is countable.
\end{corollary}
\begin{proof}
The group contains the countably infinite cyclic subgroup generated by
its least positive member and contains distinct terms $\omega^{-\beta}$
for $\beta<\alpha$. This gives the lower bound. The finite-support
sum defining $L_{\alpha,n}$ has countable coefficient groups and an
index set of size $|\alpha|+1$, giving exactly the displayed cardinality.
Its initiality and minimum follow from \cref{isg:thm:extremal}.
\end{proof}

This is a restriction on a specified initial realization. The ordinal
$\alpha$ is not an invariant of the abstract ordered group under arbitrary
isomorphisms: normalization itself changes the least positive element
to $1$.

\subsection{Dyadic depth and the spectrum of inverse scales}
\begin{definition}[Dyadic-spine depth]\label{isg:def:depth}
Let $H\subseteq\Oz$ be a nonzero initial subgroup. If there exists an
ordinal $\beta$ such that $\D\omega^{q_\beta}\not\subseteq H$, let
$d(H)$ be the least such $\beta$. Otherwise write $d(H)=\On$ as a
formal unbounded-depth marker, not as a surreal or a set ordinal.
\end{definition}

\begin{corollary}[Admissible scales]\label{isg:cor:depth}
If $d(H)$ is an ordinal, then for every $\alpha\in\On$ and every
$n\in\N$,
\[
 \Rect^{-1}[H]\text{ is initial}
 \quad\Longleftrightarrow\quad \alpha\le d(H).
\]
If $d(H)=\On$, every ordinal $\alpha$ is allowed. Whenever the inverse
is initial, its least positive member is $2^{-n}\omega^{-\alpha}$.
The parameter $n$ imposes no further obstruction.
\end{corollary}
\begin{proof}
By \cref{isg:thm:B}, the condition is exactly that all
$\D\omega^{q_\beta}$ for $\beta<\alpha$ lie in $H$. This holds
precisely through the first missing index, inclusively, and fails at all
larger indices. In the unbounded case there is no missing index.
\end{proof}

For example, $d(K_\alpha)=\alpha$, since the exponent $q_\alpha$
is absent while all the earlier required dyadic terms are present.
Thus the endpoint in \cref{isg:cor:depth} can be a limit ordinal and must
not be replaced with a strict inequality. For a set-sized $H$, a first
missing index exists: the distinct monomials $\omega^{q_\beta}$
cannot all belong to a set. The subgroup
$\bigcup_{\alpha\in\On}K_\alpha$, interpreted as the class of all
finite dyadic-spine sums with an integer constant, is an example with
unbounded depth. It is initial by \cref{isg:fact:EK}, but is much smaller
than $\Oz$, for example it does not contain $\omega^2$.

\section{The quotient by the least positive cyclic subgroup}\label{isg:sec:quotient}

\subsection{Convexity and an actual splitting}
Fix $A$, $p$, $c$, and $\varepsilon=c\omega^p$ as before. Define
\[
 \pi_p:A\longrightarrow\Z,
 \qquad \pi_p(x)=c^{-1}[\omega^p]x,
 \qquad A^{>p}=\ker\pi_p.
\]
Although its notation emphasizes the higher exponents, $A^{>p}$ also
contains zero.

\begin{lemma}[Bottom-layer decomposition]\label{isg:lem:split}
The map $\pi_p$ is a surjective additive homomorphism,
$\pi_p(m\varepsilon)=m$, and
\[
 A=A^{>p}\oplus\Z\varepsilon
\]
as additive groups. The subgroup $\Z\varepsilon$ is convex. Every
nonzero positive element of $A^{>p}$ dominates all integer multiples of
$\varepsilon$.
\end{lemma}
\begin{proof}
Coefficient extraction is additive, and $A_p=c\Z$. This proves all
assertions about $\pi_p$. Given $x\in A$, the term
$\pi_p(x)\varepsilon$ belongs to $A$ and subtraction produces an element
of $A^{>p}$. The intersection of the two displayed summands is zero.

A nonzero element of $A^{>p}$ has greatest exponent strictly above $p$.
Its absolute value therefore exceeds $m\varepsilon$ for every positive
integer $m$. If $0\le x\le m\varepsilon$ in $A$ for some integer
$m\ge0$, its component in $A^{>p}$ must consequently be zero; thus
$x\in\Z\varepsilon$. Translation and sign reversal prove convexity.
\end{proof}

The summand $A^{>p}$ need not itself be initial as a concrete subclass
of $\No$. For example, with $A=\Z+\Z\omega$, it is $\Z\omega$,
which omits the proper prefix $1$ of $\omega$. A second exponent operation
is needed for an \emph{initial} quotient realization.

\subsection{Deleting a positive root}
Every positive surreal exponent has a unique form $\plus\cat v$.
Define
\[
 \rho:\No_{>0}\longrightarrow\No,\qquad
 \rho(\plus\cat v)=v.
\]
Unlike $\tht$, this map has no exceptional point inside its domain.
It preserves comparison and preserves and reflects prefix relations
between positive source points.

\begin{lemma}[Positive-root deletion]\label{isg:lem:rho}
If $\Delta\subseteq\No_{\ge0}$ is initial, then
$E=\rho[\Delta\setminus\{0\}]$ is initial, possibly empty. For
$\delta>0$,
\[
 \Rs(\rho(\delta))=\rho[\Rs(\delta)].
\]
\end{lemma}
\begin{proof}
If $v\pr\rho(\delta)$, then $\plus\cat v\pr\delta$ is a positive
prefix and lies in $\Delta$. Its image is $v$, establishing initiality.
Every proper prefix of $\delta=\plus\cat u$ is either $0$, which is
smaller than $\delta$, or $\plus\cat v$ with $v\pr u$.
Comparisons of the latter with $\delta$ are comparisons of $v$ with $u$.
This proves the right-ancestor identity.
\end{proof}

Let $H=\Rect[A]\subseteq\Oz$ and
\[
 H^+=\{h\in H:\ct(h)=0\}.
\]
The superscript $+$ here means ``strictly positive exponents,'' not
``positive elements.'' To avoid any order-cone interpretation, define
its transported group explicitly by
\begin{equation}\label{isg:eq:Bquotient}
 B=\left\{\sum_{\delta>0}s_\delta\omega^{\rho(\delta)}:
              \sum_{\delta>0}s_\delta\omega^\delta\in H^+\right\}.
\end{equation}

\begin{proposition}[An initial model of the quotient]\label{isg:prop:quotient}
The class $B$ in \eqref{isg:eq:Bquotient} is an initial additive subgroup of
$\No$. The normal-form map in that formula induces ordered-group
isomorphisms
\[
 A/\Z\varepsilon\cong H/\Z\cong H^+\cong B.
\]
\end{proposition}
\begin{proof}
The first isomorphism is induced by $\Rect$. Since constants of $H$
are integers and all integers lie in $H$, every coset in $H/\Z$ has a
unique representative with zero constant coefficient. The quotient
ordering agrees with the ordering of those representatives, because a
nonzero positive-exponent term dominates all integers.

The group $H^+$ is truncation closed and cross-sectional on
$\Delta\setminus\{0\}$. Its coefficient groups at those points are
the same as those of $H$. Reindexing by the order-isomorphism $\rho$
produces a truncation-closed, cross-sectional subgroup with the same
coefficient groups. Its exponent class is initial by \cref{isg:lem:rho}.
The right-ancestor identity in that lemma transfers every dyadic
obligation from $H$ to $B$. Hence $B$ is initial by \cref{isg:fact:EK}.
If there are no positive exponents, the group is $\{0\}$ and the same
conclusion is immediate.
\end{proof}

\begin{proof}[Proof of \cref{isg:thm:C}]
Combine \cref{isg:lem:split,isg:prop:quotient}. The explicit isomorphism sends
$x\in A$ to
\[
 \left(\rho_*\left(\Rect(x)-\pi_p(x)\right),\ \pi_p(x)\right),
\]
where $\rho_*$ denotes the normal-form reindexing of
\eqref{isg:eq:Bquotient}. The higher component determines the sign whenever
it is nonzero; otherwise the sign is the sign of the integer component.
This is precisely the indicated lexicographic order.
\end{proof}

\begin{corollary}[Finite iteration]\label{isg:cor:iteration}
Starting with an initial group, repeatedly quotient by its least positive
cyclic subgroup whenever the current quotient is nonzero and discrete.
At every finite stage reached in this way, the quotient has an explicit
initial realization in $\No$.
\end{corollary}
\begin{proof}
Apply \cref{isg:thm:C} at the first stage and then apply it to the initial
realization just constructed at each subsequent stage. This is ordinary
finite induction and makes no assertion at a limit stage.
\end{proof}

\subsection{The converse construction: adjoining an integer bottom layer}
There is also an explicit converse to the lexicographic decomposition.
For an initial subgroup $B\subseteq\No$, including $B=\{0\}$, put
\[
 \sigma(\gamma)=\plus\cat\gamma,
 \qquad
 \sigma_*\left(\sum_\gamma r_\gamma\omega^\gamma\right)
      =\sum_\gamma r_\gamma\omega^{\sigma(\gamma)},
 \qquad J(B)=\Z\oplus\sigma_*[B].
\]
The map $\sigma_*$ changes exponents but not coefficients. Its image has
no constant term, and each nonzero image element dominates every ordinary
integer in absolute value.

\begin{theorem}[Integer-bottom suspension]\label{isg:thm:suspension}
The construction $B\mapsto J(B)$ gives an inclusion-preserving and
inclusion-reflecting correspondence between initial subgroups of $\No$
and nonzero initial subgroups of $\Oz$. Its inverse sends $H$ to the
initial group obtained by deleting the leading plus sign from every
positive exponent of $\ker(\ct|_H)$.
Moreover, $J(B)\cong B\lex\Z$ as ordered abelian groups.
\end{theorem}
\begin{proof}
For nonzero $B$, let $E=\Gamma(B)$. The exponent class of $J(B)$ is
$\{0\}\cup\sigma[E]$. It is initial: the prefixes of
$\plus\cat\gamma$ are zero and $\plus\cat v$ for $v\pr\gamma$.
The group is truncation closed and cross-sectional. Its coefficient
group at zero is $\Z$, while the coefficient group at
$\sigma(\gamma)$ is $B_\gamma$. Every right ancestor of
$\sigma(\gamma)$ has the form $\sigma(v)$ with $v\in\Rs(\gamma)$,
so its coefficient group contains $\D$. The criterion
\cref{isg:fact:EK} proves initiality. For $B=\{0\}$, the construction
is simply $J(B)=\Z$.

Conversely, if $H\subseteq\Oz$ is initial and nonzero, it contains
$\Z$ and each constant coefficient of $H$ is integral. Thus
$H=\Z\oplus\ker(\ct|_H)$. By \cref{isg:prop:quotient}, deleting the
initial plus sign on the positive exponents of that kernel gives an
initial group $B$. Reattaching the sign recovers the kernel exactly,
including its actual infinite-support membership conditions. Hence
$J(B)=H$, and the two constructions are inverse. Inclusions are
preserved and reflected because both operations are fixed reindexings
of the nonconstant normal forms. Finally,
$(b,m)\mapsto\sigma_*(b)+m$ is additive and has the stated
lexicographic order.
\end{proof}

\begin{corollary}[A relative criterion for discrete ordered groups]\label{isg:cor:relative}
Let $G$ be a nonzero discrete ordered abelian group with least positive
member $e$. Then $G$ admits an initial realization in $\No$ if and only
if $G/\Z e$ admits an initial realization and the exact sequence
\[
 0\longrightarrow\Z e\longrightarrow G
   \longrightarrow G/\Z e\longrightarrow0
\]
splits as a sequence of abelian groups.
\end{corollary}
\begin{proof}
If $G$ admits an initial realization, the least positive element is
carried to the least positive element there. Apply
\cref{isg:thm:C} and transport its splitting and quotient realization back.
For any discrete ordered group, $\Z e$ is convex: finite induction
shows that $[0,me]$ consists of $0,e,\ldots,me$ for $m\in\N$.
Conversely, a group-theoretic splitting identifies $G$ with
$(G/\Z e)\lex\Z$ as an ordered group: a lift of a nonzero positive
quotient element lies above every element of $\Z e$, by convexity.
Choose an initial realization $B$ of the quotient and apply
\cref{isg:thm:suspension}.
\end{proof}

This criterion is relative: it reduces initial realizability to the
quotient together with a splitting condition, rather than deciding
initial realizability of every possible quotient.

No result here states that every convex quotient of every initial group
has an initial realization. The cyclic least-positive case is special
because it has a normal-form coefficient projection and the exact
root-deletion operation above.

\section{Surcomplex and linear-algebraic consequences}\label{isg:sec:complex}

\subsection{Complex normal forms and rectangular modules}
Let $\No[i]=\No\oplus i\No$ with $i^2=-1$. Combining the real and
imaginary supports gives a normal form with complex coefficients and
surreal exponents; the union of the two reverse-well-ordered supports is
again reverse-well-ordered. Conjugation acts coefficientwise.
Define the Gaussian omnific module
\[
 \Oz[i]=\{a+ib:a,b\in\Oz\}.
\]
For an additive group $A\subseteq\No$, put
$A[i]=A\oplus iA$. It is a $\Z[i]$-module because multiplication by
$i$ sends $(a,b)$ to $(-b,a)$.

\begin{definition}[Rectangular initiality]\label{isg:def:rectangular}
A rectangular subclass $A+iA$ is coordinatewise initial if $A$ is
initial in $\No$. This is a deliberately coordinate-dependent notion;
no canonical single sign-tree order on $\No[i]$ is being postulated.
\end{definition}

\begin{theorem}[Surcomplex additive normalization]\label{isg:thm:complex}
Let $A$ be a nonzero discrete initial subgroup and $H=\Rect[A]$.
The map
\[
 \Rect^{\C}(a+ib)=\Rect(a)+i\Rect(b)
\]
is a $\Z[i]$-linear isomorphism
\[
 A[i]\longrightarrow H[i]\subseteq\Oz[i].
\]
It commutes with conjugation and has a complex-linear extension to the
full complex-coefficient Hahn spaces on $C_\alpha$ and
$\No_{\ge0}$. That extension preserves and reflects set-indexed strong
summability and commutes with its sums. The image is coordinatewise
initial. For fixed $\alpha,n$, its exact possible rectangular images
are the $H[i]$ with $H\subseteq\Oz$ initial and $K_\alpha\subseteq H$.
\end{theorem}
\begin{proof}
Additivity and multiplication by $i$ commute with the coordinatewise
map, so it is $\Z[i]$-linear. The inverse is coordinatewise
$\Rect^{-1}$. Conjugation negates the imaginary coordinate, and
additivity gives the asserted compatibility.

In complex normal forms, the map is exactly
\[
 \sum_{\gamma\ge p}z_\gamma\omega^\gamma
 \ \longmapsto\ c^{-1}z_p+
      \sum_{\gamma>p}z_\gamma\omega^{\tht(\gamma)}.
\]
The coefficient multiplier $c^{-1}$ is real. It commutes with every
complex scalar and with conjugation. The support argument in
\cref{isg:thm:transport} is unchanged; there is no additional infinite sum
of constants because strong summability remains point-finite.
Coordinatewise initiality follows from the initiality of $H$, and the
range assertion is exactly \cref{isg:thm:B} applied to the real factor.
\end{proof}

\subsection{What linear systems can be transported}
\begin{corollary}[Exact transport of integral linear equations]\label{isg:cor:linear}
Let $M\in\Z^{r\times s}$, let $b\in A^r$, and apply $\Rect$
coordinatewise. Then
\[
 Mx=b\quad\Longleftrightarrow\quad M\Rect(x)=\Rect(b).
\]
The corresponding statement holds for $M\in\Z[i]^{r\times s}$ and
$b\in A[i]^r$ under $\Rect^{\C}$. In the real case, coordinatewise
inequalities are preserved as well.
\end{corollary}
\begin{proof}
Every row is a finite integer or Gaussian-integer linear combination.
The relevant module isomorphism commutes with each such combination,
and its injectivity gives the reverse implication. The real order
assertion follows from the ordered-group isomorphism.
\end{proof}

The right-hand side is transported too. In general $\Rect(1)\ne1$,
so this corollary must not be rewritten as preservation of equations
with every ordinary constant fixed. Nor does it apply to nonlinear
polynomials, arbitrary surreal scalar matrices, norms, or inner products.

\section{Examples and failure modes}\label{isg:sec:examples}

\subsection{The motivating two-level example}
Ehrlich and Kaplan's example in the paragraph preceding their question
is
\[
 A=\D+\Z\omega^{-1},\qquad \varepsilon=\omega^{-1}.
\]
Our construction has $\alpha=1,n=0$, $\Theta_1(-1)=0$, and
$\Theta_1(0)=1$. It gives
\[
 d+m\omega^{-1}\longmapsto d\omega+m,
 \qquad H=\D\omega+\Z=K_1.
\]
Thus the proposed general construction specializes to their known map,
rather than treating that map as a new result \cite[Section~9]{EK}.
The example and the map stand in Section~9 of the journal version as
well, immediately before its Question~9.1.

\subsection{Rescale the bottom coefficient, not the whole series}
Consider
\[
 A=\tfrac12\Z+\Z\omega.
\]
The exponent class $\{0,1\}$ is initial, there are no right-ancestor
dyadic obligations, and its coefficient groups $\frac12\Z$ and $\Z$
are initial. Hence $A$ is initial by \cref{isg:fact:EK}. Its least positive
element is $1/2$.

Scalar multiplication by $2$ yields $\Z+2\Z\omega$, which omits
$\omega$ despite containing $2\omega$. Equivalently, its coefficient
group at exponent $1$ is not initial in $\R$. By contrast,
\[
 \mathcal N_{0,1}(a/2+b\omega)=a+b\omega
\]
has initial image $\Z+\Z\omega$. Only the bottom coefficient changes.

\subsection{An exponent translation fails at the first infinite branch}
Let
\[
 \Gamma=\{-1\}\cup\N\cup\{\omega\}
\]
be an exponent class, and let $A$ be the finite-support group with
coefficient $\Z$ at $-1$, coefficient $\D$ at $0$, and coefficient
$\Z$ at every positive member of $\Gamma$. This exponent class is
initial: $-1$ contributes the prefix $0$, and the ordinal $\omega$
contributes all finite ordinals. Its only right-ancestor obligation is
at $-1$, forcing the chosen dyadic group at zero. Thus $A$ is initial,
with least positive element $\omega^{-1}$.

Multiplication by $\omega$ replaces the support class by
\[
 \Gamma+1=\{0,1,2,\ldots\}\cup\{\omega+1\}.
\]
This class is not initial because it contains $\omega+1$ but omits
its prefix $\omega$. The resulting group is not initial. In contrast,
root relocation satisfies
\[
 \Theta_1(-1)=0,\qquad
 \Theta_1(k)=k+1\ (k<\omega),\qquad
 \Theta_1(\omega)=\omega.
\]
The last identity uses $\plus\cat\plus^\omega=\plus^\omega$;
it is precisely where ordinal concatenation differs from arithmetic
translation. The correct image has exponent class
$\N\cup\{\omega\}$, coefficient $\Z$ at zero, $\D$ at one,
and $\Z$ at the other exponents. It is initial.

\subsection{A genuine limit-support example}
Take $\alpha=\omega$, $n=0$, and consider the group of all series
\begin{equation}\label{isg:eq:limitexample}
 x=\sum_{k<\omega}a_k\omega^{-k}+m\omega^{-\omega},
 \qquad a_k\in\D,\quad m\in\Z.
\end{equation}
Zeros in this expression are ignored. The allowed support is a subset
of the reverse-well-ordered sequence
$0,-1,-2,\ldots,-\omega$. The exponent class is initial, its ancestor
coefficients are dyadic, and its bottom coefficient is integral, so the
full group \eqref{isg:eq:limitexample} is initial.

Since $\Theta_\omega(-k)=q_k=2^{-k}$ for finite $k$ and
$\Theta_\omega(-\omega)=0$, normalization gives the exact strong-sum
identity
\begin{equation}\label{isg:eq:limitimage}
 \mathcal N_{\omega,0}(x)
   =\sum_{k<\omega}a_k\omega^{2^{-k}}+m.
\end{equation}
For example, when $a_k=1$ for every $k$ and $m=1$, the source and image
have support order type $\omega+1$. The decreasing exponents
$1,1/2,1/4,\ldots,0$ are a legitimate reverse-well-ordered support;
no analytic limiting process is involved.

The finite-support minimal group $L_{\omega,0}$ is a proper subgroup of
the group in \eqref{isg:eq:limitexample}. Its image $K_\omega$ consists
only of finite sums of the positive terms appearing in
\eqref{isg:eq:limitimage}, together with an integer constant. The full
limit-support example therefore illustrates why the image cannot be
reconstructed from its coefficient groups alone.

\subsection{Inverse initiality needs the whole dyadic spine}
Let $H=\Z$ and $\alpha=1,n=0$. The inverse formula produces
\[
 \mathcal N_{1,0}^{-1}[\Z]=\Z\omega^{-1},
\]
which is not initial: $1$ is simpler than $\omega^{-1}$ but is missing.
Here $K_1\not\subseteq H$, so \cref{isg:thm:B} correctly rules out
initiality.

There is also a purely coefficient-level obstruction. Let
$H=\Z+\Z\omega$ and keep $\alpha=1,n=0$. Its exponent class contains
$\Lambda_1=\{0,1\}$, so the inverse exponent class is initial.
Nevertheless, the inverse group is
\[
 \Z+\Z\omega^{-1},
\]
which is not initial: its monomial $\omega^{-1}$ forces every dyadic
constant through \eqref{isg:eq:dyadic}, but its constant group is only
$\Z$. Equivalently, $1/2$ is a missing prefix of $\omega^{-1}$.
Thus the exponent condition in \cref{isg:prop:inversetree} alone does not
suffice for the group-level inverse.

\subsection{Multiplication and ordinary constants are not preserved}
For $A=\D+\Z\omega^{-1}$, one has $1\in A$ and
$\mathcal N_{1,0}(1)=\omega$. The product $1\cdot1$ also belongs to
$A$, but
\[
 \mathcal N_{1,0}(1\cdot1)=\omega
 \ne\omega^2=
 \mathcal N_{1,0}(1)\mathcal N_{1,0}(1).
\]
Even on products that happen to stay in the source, multiplicativity
fails. The surcomplex extension inherits this limitation. In particular,
no preservation of the complex modulus or of multiplication by general
surcomplex scalars follows from its Gaussian-linearity.

\section{Verification, provenance, and a formalization route}\label{isg:sec:audit}

\subsection{The exact dependency chain}
The proof separates into three kinds of input. Standard surreal
normal forms identify the ambient additive spaces. The established
Ehrlich--Kaplan criterion decides when a specified canonical Hahn image
is initial. Everything connecting the original discrete group to the
normalized image is then given by the explicit formulas and proofs in
this manuscript.

The essential chain is
\begin{gather*}
 \text{bottom coefficient structure}
 \ \longrightarrow\ \text{root relocation}\\
 \ \longrightarrow\ \text{strong Hahn transport}
 \ \longrightarrow\ \text{initiality of the image}.
\end{gather*}
For the inverse classification, the additional chain is
\begin{gather*}
 \text{missing ancestors of }p
 \ \longrightarrow\ \text{spine exponents}\\
 \ \longrightarrow\ \text{forced dyadic coefficients}
 \ \longrightarrow\ K_\alpha\subseteq H.
\end{gather*}
The sign-sequence proof is independent of the arithmetic of exponents.
In particular, no unproved assertion that $\tht$ preserves addition of
exponents occurs in the chain.

\subsection{Finite regression checks supplied with the article}
The accompanying Python program \texttt{verify.py} (shipped in this
report as \fn{code/verify.py}) uses only the standard
library, exact rational arithmetic, and a fixed random seed. Its recorded
run performs \textbf{450,862 assertions}, all passing (shipped as
\fn{data/verification_report.json}). The tests cover
finite sign words through length seven for all relevant finite cone
minima; every prefix-closed subset of the binary tree through depth
three, including the empty tree; inverse-spine tests on a further family
of nonnegative initial trees; and 1,500 random pairs of finite rational
normal forms.

The count of prefix-closed trees is independently specified by
\[
 T_{-1}=1,\qquad T_d=1+T_{d-1}^2,
\]
so $T_0=2$, $T_1=5$, $T_2=26$, and $T_3=677$.
The recurrence counts the empty tree, or a root with independently
chosen initial left and right subtrees. This makes the finite-tree
coverage exhaustive at the stated depth, rather than a random sample.

Tests check inverse identities, order preservation, the exact predecessor
identity, the right-ancestor
identity away from the minimum, loss of the minimum's right ancestors,
initiality of transported trees, the exact inverse exponent criterion,
and the source dyadic obligations recovered by adjoining the spine.
The series tests check inverse transport, additivity, order, truncations,
and integrality of the bottom coefficient. Separate negative tests reject
the false global right-ancestor identity and unrestricted inversion.

The program does not represent arbitrary ordinals, arbitrary real
coefficients, or infinite Hahn supports. It does not check the imported
characterization theorem or establish absence of a transfinite proof gap.
Its role is regression testing of the finite combinatorics and coefficient
formulas. The written proofs, not these experiments, carry the general
claims.

When this report was placed, the program was rerun on a copy of the
directory with an explicit output file. It passed with the same 450,862
assertions, and its JSON output was identical in content to the shipped
record. The program's default output file is
\fn{verification_report.json} in the current directory, and the
delivered build script calls it with that name, so a rerun should always
be made on a copy (the collection README gives the commands).

\subsection{Relationship to the Surreal repository}
The repository review examined the public README, the introductory
coverage and dependency ledger, and targeted searches for overlap
\cite{Repo}. Those materials describe extensive sign-sequence, canonical
cut, normal-form, Hahn-summation, and algebraic infrastructure. They also
explicitly distinguish source manuscripts from checked Lean declarations.
The present argument is designed to fit that distinction.

The review took place while the repository was changing. The initially
inspected README and ledger had blob hashes
\begin{center}\small\ttfamily
391afd7c64cebd48b47fbe43dae26c27e8ee8047\\
6a89497cb1de3e5c7a346201f7fc006d85082426.
\end{center}
A later metadata read reported commit
\begin{center}\small\ttfamily
173eb522fdf5d4bf201c097f76b38273eb4f95e1.
\end{center}
The initially returned recursive-tree hash is not used as a commit
identifier. This manuscript does not claim to have reviewed every
repository article or every companion archive, and no repository result
is silently promoted from ``pending'' to ``proved.''
The manuscript has no pinned commit. The commit \texttt{173eb52} above
is an ancestor of the commit \texttt{c6359e4} that placed the manuscript
in the collection. The manuscript makes no claim about the content of
any report, so no statement of it had to be corrected against the
collection as it now stands.

A sound implementation would first formalize the following mathematical
interfaces; the names are proposed work units, not declarations claimed
to exist:
\begin{enumerate}
\item \emph{Cone root relocation.} Define the three-way sign split,
      the inverse, the order equivalence, and
      \eqref{isg:eq:predidentity} for ordinal-length sequences.
\item \emph{Hahn reindexing with one coefficient multiplier.} Prove
      support transport, real-linearity, order preservation, truncation
      compatibility, and strong-sum interchange.
\item \emph{Concrete initial-subgroup criterion.} Supply a checked
      version of \cref{isg:fact:EK}, including both its coefficient and
      right-ancestor conditions, before claiming the final theorem.
\item \emph{Normalization and exact inverse.} Instantiate the criterion
      for the transformed groups; then derive the spine classification,
      quotient realization, and complex-coordinate extension.
\end{enumerate}

The largest imported proof obligation is the third item. A checked
conditional theorem using that criterion would be valuable, but would
not by itself constitute a checked theorem about every actual initial
subgroup of $\No$. This package supplies no Lean file with that
obligation hidden as an axiom or an unfinished proof.

\subsection{Novelty and literature status}
The new contributions claimed provisionally here are the explicit
root-relocation map together with its exact predecessor identity, its
use in omnific normalization, and the sharp dyadic-spine image
classification. The extremal, quotient, and surcomplex statements are
derived refinements of that construction. The known two-level example,
the general initial-Hahn-subgroup characterization, and the classification
of initial real coefficient groups are not claimed as new.

The literature check recovered the stated question in the original
preprint, thesis, and workshop account \cite{EK,Kaplan,Workshop}, and
examined adjacent sign-concatenation work \cite{BagHoeven}. Targeted
searches did not locate a prior resolution or the specific root-relocation
argument. That negative search result is limited evidence: it does not
exclude an unindexed paper, unpublished proof, or a reformulation already
known to a specialist. Neither mathematical correctness nor priority is
certified by the absence of a search result.

\subsection{Provenance, placement review, and place in the collection}
\label{isg:sec:provenance}

\paragraph{One source.} This report is built from one manuscript,
\emph{Discrete Initial Groups and Omnific Normalization} (24 pages,
23~September 2026), manuscript~09 of the collection's batch~28. It was
delivered with a README, a source audit, the finite checks
\texttt{verify.py}, their recorded result and a build script. It was
placed in commit \texttt{c6359e4} as this directory. The manuscript's
LaTeX source is \fn{article.tex}; its source audit is shipped unchanged as
\fn{source_audit.md}, the program and build script as \fn{code/verify.py}
and \fn{code/build.sh}, and the record as
\fn{data/verification_report.json}. The delivered PDF and README are not
shipped; the collection README replaces the latter. With one source there
was no merge and no choice between sources. The manuscript has no pinned
repository commit; its read identifiers are those recorded above.

\paragraph{What the write changed.} Every label now carries the prefix
\texttt{isg:}. The journal numbering of the Ehrlich--Kaplan results
(\cref{isg:sec:intro}) and the wording of the journal version of their
criterion (after \cref{isg:fact:EK}) were added, as were
\cref{isg:sec:conventions}, this subsection, the shipped file names, a
definition of $\ct$ (which the manuscript used without defining it),
labels on formerly unlabelled statements, and the collected list of
non-claims (\cref{isg:sec:nonclaims}). Cross-references now name each kind
of statement correctly; in the delivered source every reference to a
lemma, proposition, corollary, remark or imported fact printed as
``theorem''. No mathematical statement, proof or example was changed.

\paragraph{Placement review.} When the manuscript was placed, every step
of the argument was checked by hand: the bottom-coefficient structure,
both cases of the predecessor identity \eqref{isg:eq:predidentity} and its
failure at the old minimum, the ambient transport, every dyadic obligation
in the proof of \cref{isg:thm:A}, both directions of \cref{isg:thm:B}
with its two counterexamples, the extremal groups, the size and depth
corollaries, the splitting and suspension theorems, and every example of
\cref{isg:sec:examples}. No gap was found. The Ehrlich--Kaplan criterion
was checked as a statement, against the arXiv text and the author-hosted
text of the journal version: \cref{isg:fact:EK} is a faithful reading of
\cite[Theorem~1 and its proof]{EK}, equivalently of Theorem~5.1 of
\cite{EKjournal}, and \cref{isg:fact:real} is \cite[Lemma~3]{EK}. The
criterion was not re-proved; the whole argument rests on it. This review
is not a refereeing, and the solution proposed here has not been refereed.

\paragraph{Literature.} The placement review ran targeted web
searches (``discrete initial subgroup'', ``initial subgroup'' together
with ``omnific integers'', and the title of \cite{EK}). They found only
the Ehrlich--Kaplan paper, Kaplan's thesis and unrelated pages, and no
resolution of the question. The search was targeted, not exhaustive. The
citations of Kaplan's thesis (its Sections~6.2 and~7.1) and of the
Oberwolfach report (printed pages 3355--3356) are the manuscript's; they
were not re-checked at placement.

\paragraph{Place in the collection.} No other report names Question~2
(Question~9.1) of Ehrlich and Kaplan, and none treats initial subgroups:
at placement, a search of the reports and of the Lean sources for
``initial subgroup'', ``discrete initial'' and ``discretely ordered
initial'' found nothing outside this directory. The paper \cite{EK} is
cited as background in two reports: the foundations report cites it among
the initial-embedding results next to its sign-tree categoricity
(\replabel{found:sub:signtree}), and the definable-surreals report for
the simplicity theory of simplest separators
(\replabel{dsn:sec:foundations}). The latter's maximal initial elementary
exponential core (\replabel{dsn:thm:core}) concerns initial
\emph{subfields} of definable surreals; it is related in spirit and
neither overlaps nor conflicts with the present groups. The omnific
integers used here are those of the Diophantine and set-sized-quotient
reports (\replabel{odg:def:rings}, \replabel{osq:eq:Ozdef}); their
discreteness is \replabel{odg:prop:units}, and the constant coefficient
is the additive map of \replabel{osq:lem:ct}. The Diophantine report and
the foundations report also cite a different Ehrlich--Kaplan paper,
\emph{Surreal ordered exponential fields}, for the initial integer part
$\Oz\cap K$ of an initial subfield $K$ (the status note after
\replabel{odg:thm:floor}, and \replabel{found:sub:canonicalexp}). That
statement concerns subrings and integer parts. The groups $H\subseteq\Oz$
here are additive subgroups and $\Rect$ is not multiplicative, so there is
no conflict. No \texttt{isg:} statement has a Lean implementation.

\section{Further questions and exact boundaries}\label{isg:sec:further}

The construction establishes an additive normal form while deliberately
relinquishing several stronger kinds of compatibility. The following
questions isolate possible further work without asserting that each is
a previously published open problem.

\begin{question}[Other convex quotients]\label{isg:q:convex}
Which convex subgroups $C$ of an initial group $A\subseteq\No$ admit
an explicit initial realization of $A/C$ by a normal-form transport?
\end{question}
The proof here handles $C=\Z\varepsilon$. For a larger convex subgroup,
one would need to identify the removed exponent region and determine
whether its boundary can be relocated without creating unmet dyadic
obligations. A numerical quotient of exponent chains alone would not
settle the simplicity condition.

\begin{question}[Limit iteration]\label{isg:q:limit}
When an iterated removal of discrete bottom layers continues through
infinitely many stages, which hypotheses give a coherent initial
realization at a limit stage?
\end{question}
Finite iteration follows from \cref{isg:cor:iteration}; limit iteration
requires a separate control of supports and of the changing sign-tree
presentations. No commutation with such an unproved transfinite limit
is used in the present manuscript.

\begin{question}[Multiplicative compatibility on restricted classes]\label{isg:q:multiplicative}
For which special initial groups does some normalization to \Oz also
preserve products wherever those products are part of a specified
algebraic structure?
\end{question}
The counterexample with $1\cdot1$ rules out a general claim for our
normalization. Any useful positive result must impose additional
conditions on exponent addition and coefficient rescaling, not merely
on the order of supports.

\begin{question}[Formal certification and uniqueness principles]\label{isg:q:formal}
Can the root-relocation proof and the concrete initiality criterion be
fully verified in the repository's foundational model, and what natural
extra conditions single out this normalization among the possible
ordered-group isomorphisms to initial omnific groups?
\end{question}
Our formulas are canonical relative to the given sign model, normal
forms, and least positive scale. They are not asserted to be invariant
under arbitrary abstract group isomorphisms. The uniqueness in
\cref{isg:thm:B} concerns the inverse under the fixed map, not uniqueness of
all possible initial embeddings.

\subsection{Non-claims, collected}\label{isg:sec:nonclaims}
The limitations stated at their places in the text and in the shipped
source audit are collected here, so that none is lost. Each stays in
force where it was stated.
\begin{enumerate}[label=\textup{(N\arabic*)},leftmargin=*]
\item \emph{Status.} The manuscript is unrefereed. It is a proposed
      solution with full proposed proofs, not an independently certified
      resolution; neither correctness nor novelty is certified, and no
      priority is claimed. The placement review
      (\cref{isg:sec:provenance}) is not a refereeing.
\item \emph{Literature.} The literature and repository searches were
      targeted, not exhaustive. A negative search result is limited
      evidence and does not exclude an unindexed paper, an unpublished
      proof or a reformulation known to specialists. That a question was
      posed in 2015 and 2018 does not establish that it was still
      unresolved on the review date.
\item \emph{Imported criterion.} \Cref{isg:fact:EK} is imported, not
      re-proved, by the manuscript or at placement. No proof here
      replaces it by the weaker assertion that initiality of the exponent
      class alone suffices.
\item \emph{Not claimed as new.} The two-level example and its map,
      the initial-Hahn-subgroup characterization, the classification of
      initial real groups (\cref{isg:fact:real}) and Conway normal forms
      are imported and cited. The extremal, quotient and surcomplex
      statements are derived refinements of the construction.
\item \emph{Not a simplicity isomorphism.} $\Rect$ is an ordered-group
      isomorphism onto an initial image. It need not preserve the
      simplicity relation (\cref{isg:rem:notsimplicity}); it is not
      multiplicative, does not preserve norms, and does not fix ordinary
      constants: $\Rect(1)\ne1$ in general (\cref{isg:sec:examples}).
\item \emph{No classification.} No classification of all ordered abelian
      groups admitting initial embeddings is claimed, and the
      initial-monoid problem (Question~3 of \cite{EK}, Question~9.2 of
      the journal version) is not addressed.
\item \emph{Exponent maps.} $\Theta_\alpha$ is not a homomorphism of
      exponent groups, and \eqref{isg:eq:leadingtransport} is not
      preservation of a multiplicative valuation
      (\cref{isg:rem:preserved}). \Cref{isg:prop:length} bounds exponent
      lengths, not the birthday of a transformed normal form.
\item \emph{No internal summation.} The strong-sum statements concern the
      ambient Hahn spaces; they give no internal summation on a subgroup
      $A$ (\cref{isg:rem:nocompletion}), and \cref{isg:lem:isolate} does
      not license arbitrary infinite subseries.
\item \emph{Invariance and uniqueness.} The ordinal $\alpha$ is not an
      invariant of the abstract ordered group. The uniqueness in
      \cref{isg:thm:B} is relative to the fixed map, not uniqueness of all
      initial embeddings, and the formulas are not claimed invariant under
      abstract isomorphisms.
\item \emph{Quotients.} It is not claimed that every convex quotient of
      an initial group has an initial realization. Iteration is covered
      only at finite stages (\cref{isg:cor:iteration}), with no claim at a
      limit stage, and \cref{isg:cor:relative} is a relative criterion
      that does not decide realizability of every quotient.
\item \emph{Surcomplex scope.} \Cref{isg:thm:complex} is Gaussian-linear
      and coordinatewise initial. It gives no field isomorphism, no
      canonical sign-tree or simplicity relation on $\No[i]$, no
      preservation of the modulus, and no compatibility with
      multiplication by general surcomplex scalars. \Cref{isg:cor:linear}
      transports the right-hand side as well and does not apply to
      equations with ordinary constants held fixed, to nonlinear
      polynomials, to arbitrary surreal scalar matrices, or to norms and
      inner products.
\item \emph{Finite checks.} The finite checks do not verify transfinite
      concatenation or limit cases, arbitrary set-supported Hahn sums, the
      imported criterion, or the absence of a transfinite proof gap; the
      written proofs carry the general claims. Compiling and rendering the
      document establish its integrity, not its correctness.
\item \emph{No Lean.} No new Lean formalization and no Lean build are
      claimed. The proposed formalization units of \cref{isg:sec:audit} are
      not existing declarations, and a checked conditional theorem using
      the imported criterion would not by itself be a checked theorem
      about every initial subgroup of $\No$.
\item \emph{Repository review.} The manuscript's repository review was
      targeted: it did not review every report or archive, it treated
      truncated metadata as neither a complete inventory nor a proof of
      absence, it took no repository description as a certificate, and it
      promoted no repository result from pending to proved. The argument
      relies on no repository manuscript, and no repository Lean build was
      run for it.
\item \emph{Sources used only for context.} Bagayoko and van der Hoeven's
      sign-sequence conventions are context, not a source of the
      normalization or of its range theorem. The questions of this section
      are not asserted to be previously published open problems, and
      $d(H)=\On$ is a formal marker, not a surreal number or a set
      ordinal.
\end{enumerate}

\section{Conclusion}\label{isg:sec:conclusion}
The obstruction in the Ehrlich--Kaplan question is not the numerical
size of the least positive element. Scalar multiplication removes that
size immediately. The obstruction is the interaction of the exponent
tree with coefficient groups forced by right ancestors.

The explicit map $\Theta_\alpha$ moves the minimum exponent to zero
and preserves all right-ancestor relations away from it. The coefficient
map then normalizes the discrete bottom group without disturbing the
higher groups. This gives the proposed affirmative proof. The relations
lost at the minimum are exactly those recorded by the dyadic spine,
which yields a necessary-and-sufficient image criterion rather than
only a sufficient construction.

The resulting theory is additive but strong: it transports arbitrary
set-supported normal forms without changing their support order types,
commutes with strong sums, identifies extremal initial groups at each
least positive scale, and produces initial models of the bottom cyclic
quotients and rectangular surcomplex modules. The remaining obligations
are independent mathematical review, a more exhaustive priority check,
and a complete formal verification of the identified dependency chain.

\appendix
\section{A compact proof for independent checking}\label{isg:app:compact}

This appendix compresses the main existence argument without changing
its hypotheses. The full case analysis and support arguments are in the
body of the paper.

Let $A$ be a discrete initial subgroup. Its least positive member has a
single normal-form term: otherwise a nonzero tail in $A$ would be a
smaller positive element after taking absolute value. Write that term
as $c\omega^p$. The exponent class has minimum $p$, which, being the
minimum of an initial tree, is all-minus, say $p=-\alpha$.
The bottom real coefficient group is discrete and initial, hence
$c\Z$ with $c=2^{-n}$.

On $[p,\infty)$ define
\[
 p\mapsto0,\qquad
 p\cat\plus\cat u\mapsto\plus\cat p\cat u,\qquad
 x\mapsto\plus\cat x\quad(p\not\pre x).
\]
The inverse removes the initial plus and reinserts a plus immediately
after the block $p$ whenever the remaining word starts with $p$.
The two maps are increasing inverse bijections. For $x>p$,
\[
 \Pred(\Theta_\alpha x)
   =\Theta_\alpha[\Pred(x)\cup\{p\}],
 \qquad
 \Rs(\Theta_\alpha x)=\Theta_\alpha[\Rs(x)].
\]
Therefore an initial exponent class with minimum $p$ maps to an initial
one with minimum zero, and every new right-ancestor obligation comes
from an old one above $p$.

Transport each nonbottom term $r\omega^\gamma$ to
$r\omega^{\Theta_\alpha\gamma}$, and transport the bottom term
$r\omega^p$ to $r/c$. This is a real-linear order-isomorphism on the
ambient Hahn spaces. Its image of $A$ is truncation closed and
cross-sectional, has coefficient group $\Z$ at zero, and has unchanged
initial real coefficient groups at positive exponents. The right-ancestor
identity preserves all required dyadic containments. By the concrete
Ehrlich--Kaplan criterion, the image is initial. It lies in \Oz because
all exponents are nonnegative and the constant coefficient is integral.

For the inverse classification, only the ancestors of the old minimum
were discarded. They are $-\beta$ for $\beta<\alpha$, and their images
are $q_\beta=\plus\cat\minus^\beta$. Initial inversion requires
those exponents, while the right-ancestor condition at $p$ requires all
their dyadic coefficients. Together these requirements are exactly
$K_\alpha\subseteq H$.

\clearpage
\section{Proof-review checklist}\label{isg:app:checklist}
An independent review can focus on the following specific potential
failure points rather than checking experimental output as though it
were a proof.

\begin{longtable}{@{}p{0.30\textwidth}p{0.64\textwidth}@{}}
\toprule
Issue & Location and required check\\
\midrule
\endhead
Concrete versus existential initiality &
\Cref{isg:fact:EK}: the imported sufficiency statement must apply to the
specified canonical image, not only to an unspecified isomorphic copy.\\[4pt]
Minimum exponent &
\Cref{isg:prop:bottom}: truncation and subtraction isolate a nonzero tail;
its absolute value is smaller than the positive leading monomial.\\[4pt]
Limit sign lengths &
\Cref{isg:lem:inverse,isg:thm:tree}: concatenated tails use ordinal interval
order types. No subtraction of ordinal lengths or removal of a last
limit sign is permitted.\\[4pt]
Exceptional old minimum &
\Cref{isg:thm:tree,isg:rem:exception}: the right-ancestor equality excludes
$x=p$. Its failure is used, not ignored, in the inverse classification.\\[4pt]
Coefficient group scaling &
\Cref{isg:lem:coefficients}: only the bottom coefficient is rescaled;
all right-ancestor obligations in the target are at unchanged groups.\\[4pt]
Infinite subgroup membership &
\Cref{isg:sec:transport}: transport the actual subgroup. Never replace it
with all normal forms allowed by its coefficient groups.\\[4pt]
Sharp inverse criterion &
\Cref{isg:thm:B}: inverse exponent initiality is not enough; the missing
minimum's ancestors must carry all dyadic coefficients.\\[4pt]
Quotient order &
\Cref{isg:lem:split,isg:prop:quotient}: higher nonzero terms dominate the cyclic
bottom group, and positive-root deletion supplies an initial image.\\[4pt]
Class-sized scope &
All supports and summation indices are sets; correspondences of
class-sized groups are expressed with class parameters.\\[4pt]
Surcomplex scope &
\Cref{isg:thm:complex}: the conclusion is Gaussian-linear and coordinatewise
initial, not a field isomorphism or a new canonical complex simplicity
relation.\\
\bottomrule
\end{longtable}

\clearpage
\phantomsection
\addcontentsline{toc}{section}{References}
\begin{thebibliography}{9}
\bibitem{EK}
P.~Ehrlich and E.~Kaplan,
\emph{Number systems with simplicity hierarchies: a generalization of
Conway's theory of surreal numbers II},
arXiv:1512.04001v1, 13 December 2015.
\url{https://arxiv.org/abs/1512.04001v1}.
The version explicitly used for theorem numbering: Lemma~3,
Theorem~1 and its concrete sufficiency proof, Proposition~9,
and Section~9, Question~2.

\bibitem{EKjournal}
P.~Ehrlich and E.~Kaplan,
\emph{Number systems with simplicity hierarchies: a generalization of
Conway's theory of surreal numbers II},
The Journal of Symbolic Logic \textbf{83} (2018), no.~2, 617--633.
\url{https://doi.org/10.1017/jsl.2017.9}.
Publisher metadata checked; the arXiv version above supplies the cited
internal numbering. The journal numbering given in \cref{isg:sec:intro}
(Lemma~5.1, Theorem~5.1, Proposition~5.2, Questions~9.1 and~9.2) was
read, at placement, in the authors' text of this version hosted on the
second author's web page; the publisher's typeset pages were not
compared.

\bibitem{Kaplan}
E.~Kaplan,
\emph{Initial Embeddings in the Surreal Number Tree},
Honors thesis, Ohio University, 2015.
OhioLINK accession \texttt{ouhonors1429615758}.
See Sections~6.2 and~7.1 for omnific integers and the group question.
\url{https://etd.ohiolink.edu/acprod/odb_etd/ws/send_file/send?accession=ouhonors1429615758&disposition=inline}.

\bibitem{Workshop}
\emph{Mini-Workshop: Surreal Numbers, Surreal Analysis, Hahn Fields
and Derivations}, Oberwolfach Report 60/2016,
workshop account containing the initial-subgroup question and the
concrete Hahn criterion, especially printed pages 3355--3356.
\url{https://ems.press/content/serial-article-files/46663}.

\bibitem{BagHoeven}
V.~Bagayoko and J.~van der Hoeven,
\emph{Surreal substructures}, author-hosted manuscript,
first version 8 June 2019, corrected version 16 April 2022.
\url{https://www.texmacs.org/joris/sss/sss.html}.
See the sign-sequence conventions and concatenation operations.

\bibitem{Repo}
V.~Reshetnikov,
\emph{Surreal}, public research and Lean-formalization repository,
reviewed 23 September 2026.
\url{https://github.com/VladimirReshetnikov/Surreal}.
Reviewed entry points: \texttt{README.md} and
\texttt{docs/FORMALIZATION.md}; precise read identifiers and the scope
of the review are recorded in \cref{isg:sec:audit} and the accompanying
\texttt{source\_audit.md}.
\end{thebibliography}
\end{document}
